\documentclass[12pt]{article}

\usepackage[T1]{fontenc}
\usepackage[margin=1in]{geometry}
\usepackage{setspace}
\usepackage{parskip}
\usepackage{amsmath}
\setstretch{1.0}

\title{Lab 12}
\author{Chuan Wei Heng}
\date{\today}

\begin{document}
\maketitle

\section*{3.2 Exercises}

We compare \texttt{scipy.linalg.solve} and LU factorization for the square
system $Ax=b$, using $N=100,500,1000,2000,4000,$ and $5000$. In an
initial test with $N=5000$, the direct solver took $1.1674$ s and the
combined LU method took $0.7059$ s. We still need the $N=4000$ result
and separate timings for LU factorization and triangular substitution.

For $k$ different right-hand sides, reusing the LU factorization is
faster when
\[
k>\frac{t_{\mathrm{factor}}}
{t_{\mathrm{direct}}-t_{\mathrm{triangular}}}.
\]
The numerical break-even count depends on the separate timings.

\section*{3.4 Exercises}

For an overdetermined system ($N>M$), we compare the normal equations
with reduced QR factorization:
\[
A^TAx=A^Tb,\qquad A=QR,\quad Rx=Q^Tb.
\]
We validate both by choosing $x_{\mathrm{true}}$, setting
$b=Ax_{\mathrm{true}}$, and checking the residual and solution error.

With the original smallest singular value $10^{-10}$,
$\kappa_2(A)=10^9$ and $\kappa_2(A^TA)=10^{18}$. Reducing it to
$10^{-12}$ increases these to $10^{11}$ and $10^{22}$, respectively.
QR should be more accurate because the normal equations square the
condition number. Numerical errors still need to be recorded.

\section*{Conclusion}

Reusing LU is advantageous for multiple right-hand sides, while QR is
more reliable for ill-conditioned least-squares problems. The completed
code must be pushed to the \textit{chuanweiheng-cell} APPM4600 Git
repository, and the answers submitted on Canvas.

\end{document}
